%==============================================================================
% Chapitre 71 - L'accélération du projectile dans le canon
%==============================================================================

\section{Introduction}

Une fois le projectile mis en mouvement, les gaz issus de la combustion
continuent à exercer une pression sur son culot.

Cette pression lui communique une accélération qui se poursuit pendant la plus
grande partie de son parcours dans le canon.

\begin{aretenir}

Le projectile accélère tant que la force exercée par les gaz demeure
supérieure aux forces qui s'opposent à son mouvement.

\end{aretenir}

\section{Origine de l'accélération}

La pression des gaz exerce une force dirigée suivant l'axe du canon.

Cette force est responsable de l'accélération du projectile.

À mesure que le projectile avance, sa vitesse augmente continuellement jusqu'à
sa sortie de bouche.

\begin{figure}[htbp]
\centering

\begin{tikzpicture}[>=Latex]

% Projectile
\fill[ProjectileColor]
(0,-0.4)
rectangle
(3,0.4);

\node at (1.5,0){Projectile};

% Gaz
\fill[GasColor]
(-3,-0.4)
rectangle
(0,0.4);

\node at (-1.5,0){Gaz};

% Force pression
\draw[-{Latex[length=4mm]},very thick,PressureColor]
(-1,0.6)--(0.8,0.6);

\node[above]
at (-0.1,0.6)
{$F_{\mathrm{pression}}$};

% Frottement
\draw[-{Latex[length=4mm]},very thick]
(4,0)--(3,0);

\node[right]
at (4,0)
{$F_{\mathrm{frottement}}$};

\end{tikzpicture}

\caption{Forces principales agissant sur le projectile.}
\label{fig:forces_projectile}

\end{figure}

\begin{formule}{Force exercée par les gaz}

La force appliquée au culot du projectile est égale au produit de la pression
par la surface sur laquelle elle s'exerce.

\begin{description}
\item[$F$] Force exercée sur le projectile (N)
\item[$P$] Pression des gaz (Pa)
\item[$S$] Surface du culot ($m^2$)
\end{description}

\end{formule}

\section{Deuxième loi de Newton}

Le mouvement du projectile est régi par les lois de la mécanique classique.

La force résultante appliquée au projectile détermine directement son
accélération.

\begin{rigueur}{\textbf{Force et accélération}}

Une même force ne produit pas la même accélération sur tous les projectiles.

À force identique, un projectile plus massif accélère moins qu'un projectile
plus léger.

Les relations mathématiques correspondantes seront développées dans un chapitre
ultérieur.

\end{rigueur}

\begin{figure}[htbp]
\centering

\begin{tikzpicture}[
block/.style={
draw,
rounded corners,
minimum width=32mm,
minimum height=8mm
}
]

\node[block](P){Pression};

\node[block,right=22mm of P](F){Force};

\node[block,right=22mm of F](A){Accélération};

\draw[-Latex,thick](P)--(F);
\draw[-Latex,thick](F)--(A);

\end{tikzpicture}

\caption{La pression engendre une force qui provoque l'accélération du projectile.}

\end{figure}

\begin{formule}{Deuxième loi de Newton}

L'accélération est proportionnelle à la force résultante appliquée au
projectile et inversement proportionnelle à sa masse.

\begin{equation}
\label{newton2}
F = ma
\end{equation}

\begin{description}

\item[$F$] Force résultante (N)

\item[$m$] Masse du projectile (kg)

\item[$a$] Accélération ($m.s^{-2}$)

\end{description}

\end{formule}

\section{Évolution de la vitesse}

La vitesse n'augmente pas de manière parfaitement linéaire.

Elle dépend notamment :

\begin{itemize}
    \item de la pression des gaz ;
    \item de la combustion de la poudre ;
    \item de la masse du projectile ;
    \item des frottements.
\end{itemize}

\begin{figure}[htbp]
\centering

\begin{tikzpicture}

\draw[-Latex]
(0,0)--(10,0)
node[right]{Temps};

\draw[-Latex]
(0,0)--(0,5)
node[above]{Vitesse};

\draw[very thick,ProjectileColor]
plot[smooth]
coordinates{
(0.3,0.1)
(1,0.4)
(2,1.2)
(3,2.3)
(4,3.3)
(5,4)
(6,4.4)
(7,4.7)
(8,4.9)
(9,5)
};

\end{tikzpicture}

\caption{Évolution qualitative de la vitesse du projectile.}

\end{figure}

\section{Influence des frottements}

Le déplacement du projectile s'accompagne de frottements entre celui-ci et
l'intérieur du canon.

Ces frottements absorbent une partie de l'énergie disponible.

Ils contribuent également à l'échauffement du canon et du projectile.

\section{Influence des rayures}

Dans un canon rayé, une partie de l'énergie est utilisée pour communiquer un
mouvement de rotation au projectile.

Cette énergie n'est donc plus disponible pour son accélération longitudinale.

Cependant, cette perte reste relativement faible devant l'énergie totale mise
en jeu.

\section{Évolution de la pression}

Pendant que le projectile accélère, deux phénomènes évoluent simultanément :

\begin{itemize}
    \item la combustion continue à produire des gaz ;
    \item le volume situé derrière le projectile augmente.
\end{itemize}

La pression résulte en permanence de l'équilibre entre ces deux phénomènes.

\section{Accélération maximale}

L'accélération du projectile peut atteindre des valeurs extrêmement élevées.

Elle dépasse largement celles rencontrées dans la plupart des applications
courantes de la mécanique.

Cette phase ne dure toutefois que quelques millisecondes.

\section{Fin de l'accélération}

Lorsque le projectile approche de la bouche du canon, la pression diminue
progressivement.

L'accélération décroît alors jusqu'à devenir pratiquement nulle au moment où le
projectile quitte le canon.

\section{Observation expérimentale}

L'accélération peut être déterminée indirectement à partir de mesures de
vitesse réalisées le long du canon ou à sa sortie.

Les résultats sont utilisés pour valider les modèles de balistique intérieure.

\section{Représentation en simulation}

\begin{simulation}

Les simulateurs calculent à chaque pas de temps la force exercée sur le
projectile, puis en déduisent son accélération, sa vitesse et sa position.

Cette boucle de calcul est répétée jusqu'à la sortie du projectile.

\end{simulation}

\section{Vers la rotation du projectile}

Dans un canon rayé, le déplacement longitudinal s'accompagne d'une mise en
rotation imposée par les rayures.

Cette rotation jouera un rôle essentiel dans la stabilité du projectile après
sa sortie du canon.

\section{À retenir}

\begin{aretenir}

\begin{itemize}
    \item Les gaz accélèrent le projectile tout au long de son parcours.
    \item L'accélération dépend de la force résultante exercée sur le projectile.
    \item Les frottements et les rayures consomment une partie de l'énergie.
    \item La vitesse augmente jusqu'à la sortie du canon.
    \item La rotation est acquise simultanément au mouvement longitudinal.
\end{itemize}

\end{aretenir}
